\chapter{Электромагнитные волны и действие Максвелла}

\lectureoverview{Сначала мы решим уравнения Максвелла без источников для плоской волны и позволим самим уравнениям определить ее поляризацию, геометрию поля и скорость. Затем мы зададим более глубокий вопрос: какое локальное, лоренц-инвариантное, калибровочно-инвариантное действие порождает динамические уравнения Максвелла? Ответ связывает электромагнитное поле, сохранение заряда и калибровочную инвариантность в единую структуру.}

\section{Два семейства уравнений Максвелла}

Начнем с плоской электромагнитной волны. Здесь \emph{plane} — геометрическое слово, а не \emph{plain} в смысле простой, без украшений. Непосредственные вопросы здесь вполне конкретны: как выглядит такая волна и как ее форма следует из уравнений Максвелла?

Полезно, прежде чем делать с волной что-либо еще, выписать уравнения в два столбца. В пустом пространстве они имеют вид
\begin{align}
\nabla\cdot\mathbf B&=0,
&
\dot{\mathbf B}&=\nabla\times\mathbf E,
\label{eq:homogeneous-vector}
\\
\nabla\cdot\mathbf E&=0,
&
\dot{\mathbf E}&=-\,\nabla\times\mathbf B.
\label{eq:dynamical-vector-vacuum}
\end{align}
Мы используем метрику с преимущественно положительной сигнатурой
\begin{equation}
\eta_{\mu\nu}=\operatorname{diag}(-1,1,1,1)
\end{equation}
и принятое в курсе соглашение
\begin{equation}
F_{\mu\nu}
=
\partial_\mu A_\nu-\partial_\nu A_\mu,
\qquad
F_{0i}=E_i,
\qquad
F_{ij}=\epsilon_{ijk}B_k.
\label{eq:course-f-convention}
\end{equation}
Тогда те же два семейства можно записать как
\begin{align}
\partial_\mu F_{\nu\sigma}
+\partial_\nu F_{\sigma\mu}
+\partial_\sigma F_{\mu\nu}
&=0,
\label{eq:bianchi-l9}
\\
\partial_\mu F^{\mu\nu}
&=0.
\label{eq:vacuum-maxwell-covariant}
\end{align}

\begin{figure}[H]
\centering
\includegraphics[width=\linewidth]{lecture_09_figure_01.png}
\caption{Два вакуумных семейства уравнений Максвелла и их ковариантные формы. Циклическое уравнение является тождеством; свернутое уравнение — динамическим.}
\par\smallskip
\textit{Кадр лекции: 00:03:30.}
\end{figure}

Различие между двумя строками существенно. Уравнение
\eqref{eq:bianchi-l9} тождественно следует из существования потенциала \(A_\mu\): каждый член является второй производной, и эти члены попарно сокращаются. Следовательно, это тождество Бьянки. Уравнение
\eqref{eq:vacuum-maxwell-covariant}, напротив, является уравнением движения. Измените электромагнитный лагранжиан — и изменится именно это семейство.

На доске есть небольшое наблюдение, которое в этой главе выполнит две задачи. Каждое поле в уравнениях Максвелла появляется под знаком производной. Этот факт наложит ограничения на относительные фазы плоской волны и позднее сделает естественной калибровочную инвариантность.

\section{Плоская волна вдоль оси \texorpdfstring{\(z\)}{z}}

Выберем координаты так, чтобы волна распространялась вдоль оси \(z\). Это не ограничивает общности; волну, идущую в другом направлении, можно привести к такому виду поворотом осей. Пусть \(k\) — ее волновое число,
\begin{equation}
\lambda=\frac{2\pi}{k},
\end{equation}
а \(\omega\) — ее частота. Прежде чем налагать какое-либо уравнение движения, запишем все шесть компонент поля с общей фазой
\begin{equation}
\theta=kz-\omega t:
\end{equation}
\begin{align}
E_x(z,t)&=\epsilon_x\sin\theta,
&
B_x(z,t)&=\beta_x\sin\theta,
\nonumber\\
E_y(z,t)&=\epsilon_y\sin\theta,
&
B_y(z,t)&=\beta_y\sin\theta,
\label{eq:plane-wave-ansatz}\\
E_z(z,t)&=\epsilon_z\sin\theta,
&
B_z(z,t)&=\beta_z\sin\theta.
\nonumber
\end{align}
Эти шесть греческих букв — постоянные амплитуды.

\begin{figure}[H]
\centering
\includegraphics[width=\linewidth]{lecture_09_figure_02.png}
\caption{Анзац плоской волны с шестью компонентами до того, как уравнения Максвелла исключат недопустимые амплитуды.}
\par\smallskip
\textit{Кадр лекции: 00:08:45.}
\end{figure}

Почему мы начинаем с одной и той же синусоидальной зависимости для \(\mathbf E\) и
\(\mathbf B\)? Предположим, что одно поле содержит \(\sin\theta\), а другое —
\(\cos\theta\). Тогда производная по времени с одной стороны уравнения ротора и пространственная производная с другой дали бы разные тригонометрические функции. Равенство не могло бы выполняться при всех \(z\)
и \(t\). Общая постоянная фаза допустима, но сдвиг начала отсчета
\(z\) ее устраняет. Таким образом, этот анзац — не дополнительный динамический
закон; это простейший представитель того фазового соотношения, которое
потребуют дифференциальные уравнения.

Пока ни одно уравнение Максвелла не использовано. Мы лишь выбрали направление,
длину волны и удобное начало отсчета фазы. Теперь уравнения решат,
какие амплитуды сохраняются.

\subsection{Ограничения на дивергенцию}

Поскольку поля зависят только от \(z\) и \(t\),
\begin{equation}
\nabla\cdot\mathbf E
=
\partial_xE_x+\partial_yE_y+\partial_zE_z
=
\partial_zE_z.
\end{equation}
В пустом пространстве это должно обращаться в нуль при любой фазе, так что
\begin{equation}
\epsilon_z=0.
\label{eq:epsilon-z-zero}
\end{equation}
Тот же самый довод для магнитного поля дает
\begin{equation}
\beta_z=0.
\label{eq:beta-z-zero}
\end{equation}
Ни одно из полей не имеет компоненты вдоль направления распространения. Волна поперечна.

В поперечной плоскости \(xy\) остается свобода поворота. Пара \((\epsilon_x,\epsilon_y)\) — обычный двумерный вектор,
поэтому повернем поперечные оси так, чтобы
\begin{equation}
\epsilon_y=0,
\qquad
\epsilon_x\ne0.
\label{eq:polarization-axis}
\end{equation}
Это выбор линейной поляризации, а не ограничение физического содержания задачи.

\begin{figure}[H]
\centering
\includegraphics[width=\linewidth]{lecture_09_figure_03.png}
\caption{Продольные амплитуды исчезли, и выбрана поперечная электрическая ось.}
\par\smallskip
\textit{Кадр лекции: 00:14:30.}
\end{figure}

\subsection{Уравнения ротора}

Возьмем \(x\)-компоненту уравнения
\(\dot{\mathbf E}=-\nabla\times\mathbf B\). Поскольку поля зависят только
от \(z\),
\begin{equation}
\dot E_x
=
-\left(\nabla\times\mathbf B\right)_x
=
\partial_zB_y.
\end{equation}
Подстановка анзаца дает
\begin{align}
-\omega\epsilon_x\cos\theta
&=
k\beta_y\cos\theta,
\nonumber\\
\beta_y
&=
-\,\epsilon_x\frac{\omega}{k}.
\label{eq:beta-y-epsilon}
\end{align}
\(y\)-компонента того же уравнения Максвелла имеет вид
\begin{equation}
\dot E_y
=
-\left(\partial_zB_x\right).
\end{equation}
Левая часть равна нулю в силу нашего выбора поляризации, так что
\begin{equation}
\beta_x=0.
\label{eq:beta-x-zero}
\end{equation}

\begin{figure}[H]
\centering
\includegraphics[width=\linewidth]{lecture_09_figure_04.png}
\caption{Покомпонентный расчет ротора связывает \(\epsilon_x\) с
\(\beta_y\) и устраняет оставшуюся магнитную амплитуду.}
\par\smallskip
\textit{Кадр лекции: 00:20:15.}
\end{figure}

Геометрия тем самым фиксирована:
\begin{equation}
\mathbf E\parallel\hat{\mathbf x},
\qquad
\mathbf B\parallel\hat{\mathbf y},
\qquad
\mathbf k\parallel\hat{\mathbf z}.
\end{equation}
Электрическое поле, магнитное поле и направление распространения взаимно перпендикулярны.

\begin{figure}[H]
\centering
\includegraphics[width=0.94\linewidth]{lecture_09_figure_05.png}
\caption{Изображение на доске поперечных электрических и магнитных колебаний, распространяющихся вдоль оси \(z\).}
\par\smallskip
\textit{Кадр лекции: 00:23:15.}
\end{figure}

Другое уравнение ротора дает независимое соотношение между амплитудами. Его \(y\)-компонента имеет вид
\begin{equation}
\dot B_y=(\nabla\times\mathbf E)_y=\partial_zE_x,
\end{equation}
и, следовательно,
\begin{align}
-\omega\beta_y\cos\theta
&=
k\epsilon_x\cos\theta,
\nonumber\\
\epsilon_x
&=
-\,\beta_y\frac{\omega}{k}.
\label{eq:epsilon-beta-y}
\end{align}

\begin{figure}[H]
\centering
\includegraphics[width=\linewidth]{lecture_09_figure_06.png}
\caption{Взаимные соотношения амплитуд над эскизом поперечной волны.}
\par\smallskip
\textit{Кадр лекции: 00:25:00.}
\end{figure}

Объединяя уравнения \eqref{eq:beta-y-epsilon} и
\eqref{eq:epsilon-beta-y}, получаем дисперсионное соотношение
\begin{equation}
\left(\frac{\omega}{k}\right)^2=1.
\end{equation}
Для волны с положительной частотой, движущейся в выбранном направлении,
\begin{equation}
\omega=k
\qquad
(c=1).
\label{eq:omega-k}
\end{equation}
Поэтому общая фаза равна
\begin{equation}
\theta=k(z-t).
\end{equation}
Возвращаясь к обычным единицам, получаем
\begin{equation}
\omega=ck,
\qquad
\theta=k(z-ct).
\label{eq:restore-c-wave}
\end{equation}
Тем самым профиль переносится без искажения со скоростью света. В единицах, использованных на доске,
электрическая и магнитная амплитуды имеют одинаковую величину.
Их относительный знак минус — это соглашение об ориентации: он указывает,
в какую сторону направлена выбранная ось \(y\), а не на то, что одно поле отстает от другого на четверть периода.

\begin{figure}[H]
\centering
\includegraphics[width=\linewidth]{lecture_09_figure_07.png}
\caption{Окончательный вид плоской волны после наложения условия \(\omega=k\).}
\par\smallskip
\textit{Кадр лекции: 00:26:45.}
\end{figure}

\begin{classroomqa}
\audiencequestion{Разве размерности \(\omega\) и \(k\) не различны?}
\lecturerresponse{Различны в обычных единицах: \(\omega\) имеет размерность обратного времени,
а \(k\) — обратной длины. При \(c=1\) время и длина имеют одну и ту же
единицу. Восстановление \(c\) превращает \(\omega=k\) в \(\omega=ck\), в точности как и
требует размерностный анализ. \lecturetimestamp{00:25:55}}
\end{classroomqa}

\begin{classroomqa}
\audiencequestion{Всегда ли электрическое и магнитное поля в плоской волне находятся в одной фазе?}
\lecturerresponse{Их фазовая зависимость одинакова. Их знаковые
амплитуды могут отличаться знаком минус, что можно описать как
разворот на \(180^\circ\), но уравнения Максвелла не допускают
\(90^\circ\) сдвига синус--косинус для этой линейно поляризованной плоской волны.
Круговая поляризация составляется из поперечных компонент и
требует более аккуратного языка, но она не изменяет этого покомпонентного
соотношения Максвелла. \lecturetimestamp{00:27:31}}
\end{classroomqa}

\begin{classroomqa}
\audiencequestion{Если оба поля одновременно проходят через нуль, сохраняется ли
энергия?}
\lecturerresponse{Да. Плотность энергии в одной фиксированной точке осциллирует,
когда волна проходит мимо, но энергия переместилась, а не исчезла.
То, что в один момент времени занимало одну область, мгновение спустя занимает соседнюю.
Локальный перенос энергии, более полно представляемый потоком
Пойнтинга, согласует изменяющуюся плотность с законом сохранения.
\lecturetimestamp{00:28:15}}
\end{classroomqa}

Уравнения Максвелла линейны. Амплитуда \(\epsilon_x\) может быть
большой или малой, \(\beta_y\) из нее следует, и любая сумма или разность
решений снова является решением. Поэтому волны, распространяющиеся в разных направлениях,
могут накладываться друг на друга. Этот короткий расчет стоит проделать:
изложение электромагнетизма не должно оставлять электромагнитную волну
просто необъясненной цитатой.

\begin{classroomqa}
\audiencequestion{Предполагалось ли в выводе с самого начала, что
\(\mathbf E\) и \(\mathbf B\) имеют одну и ту же фазу?}
\lecturerresponse{Общая фаза была записана как анзац, но она
не является независимым физическим предположением. Поскольку на каждой стороне уравнений ротора
берется производная от поля, уравнения исключают нетривиальную
относительную фазу. После этого общий фазовый сдвиг можно устранить сдвигом начала отсчета. \lecturetimestamp{00:29:56}}
\end{classroomqa}

\begin{classroomqa}
\audiencequestion{Откуда взялось второе соотношение
\(\epsilon_x=-\beta_y\omega/k\)?}
\lecturerresponse{Первое соотношение следовало из
\(\dot{\mathbf E}=-\nabla\times\mathbf B\). Второе получается из
сопутствующего уравнения \(\dot{\mathbf B}=\nabla\times\mathbf E\), где
\(\mathbf E\) и \(\mathbf B\) меняются местами в том же покомпонентном
вычислении. Для получения \(\omega^2=k^2\) нужны оба.
\lecturetimestamp{00:30:25}}
\end{classroomqa}

\section{Почему дифференциальных уравнений недостаточно}

Теперь мы переходим от частного решения к структуре теории.
Именно принцип наименьшего действия организует закон сохранения энергии,
закон сохранения импульса и их связь с симметриями пространства-времени.
Можно придумать дифференциальные уравнения, у которых нет энергии, которую стоило бы
назвать сохраняющейся, и нет никакого полезного принципа симметрии. Уравнения Максвелла,
напротив, должны следовать из лагранжиана.

Нужно вывести только динамическое семейство:
\begin{equation}
\partial_\mu F^{\mu\nu}=J^\nu.
\label{eq:target-sourced-maxwell}
\end{equation}
Семейство Бианки уже следует из определения \(F_{\mu\nu}\).
На языке трёхвекторов наша конвенция даёт
\begin{align}
\nabla\cdot\mathbf E&=\rho,
\label{eq:gauss-source}
\\
\dot{\mathbf E}+\nabla\times\mathbf B&=-\,\mathbf J.
\label{eq:ampere-course-sign}
\end{align}
Знак минус в \eqref{eq:ampere-course-sign} важен. Это
покомпонентная форма \eqref{eq:target-sourced-maxwell} при
\eqref{eq:course-f-convention}, и мы снова обнаружим его при непосредственной
проверке знака ближе к концу.

\begin{figure}[H]
\centering
\includegraphics[width=\linewidth]{lecture_09_figure_08.png}
\caption{Полная запись уравнений Максвелла на доске при переходе от волнового решения к принципу действия.}
\par\smallskip
\textit{Кадр лекции: 00:32:45.}
\end{figure}

\section{Локальная релятивистская теория поля}

Первый принцип — локальность. Действие есть интеграл по пространству-времени
от плотности, построенной из полей и их первых производных в той же
точке:
\begin{equation}
S=\int d^4x\,\mathcal L(\phi,\phi_{,\mu}),
\qquad
\phi_{,\mu}\equiv\partial_\mu\phi.
\label{eq:local-action}
\end{equation}
Запятая — это компактное обозначение производной. Второй принцип —
лоренц-инвариантность: \(\mathcal L\) должен быть скаляром.

\subsection{Тренировка на скалярном поле}

Прежде чем варьировать четырёхвектор, отработаем метод на одном скалярном поле.
Локальный лоренц-скаляр может зависеть от потенциала \(U(\phi)\) и от
свёртки двух производных:
\begin{equation}
\mathcal L_\phi
=
-\frac12\,\partial_\mu\phi\,\partial^\mu\phi-U(\phi).
\label{eq:scalar-lagrangian}
\end{equation}
При метрике с тремя плюсами,
\begin{equation}
\partial_\mu\phi\,\partial^\mu\phi
=
-\dot\phi^{\,2}+|\nabla\phi|^2,
\end{equation}
поэтому
\begin{equation}
\mathcal L_\phi
=
\frac12\dot\phi^{\,2}
-\frac12|\nabla\phi|^2
-U(\phi).
\end{equation}
Производная по времени с поднятым индексом отличается по знаку от производной с опущенным; производные по пространственным координатам с поднятыми
и опущенными индексами не отличаются.

\begin{classroomqa}
\audiencequestion{Имеет ли какой-нибудь смысл писать индекс производной
сверху, а не снизу?}
\lecturerresponse{Да. Поднятие индекса означает свёртку с обратной
метрикой. При сигнатуре \((-+++)\),
\(\partial^0=-\partial_0\), тогда как
\(\partial^i=\partial_i\). Именно поэтому член с производной
по времени несёт знак, противоположный трём пространственным членам.
\lecturetimestamp{00:38:02}}
\end{classroomqa}

Уравнение Эйлера--Лагранжа для поля имеет вид
\begin{equation}
\partial_\mu\left(
\frac{\partial\mathcal L}{\partial\phi_{,\mu}}
\right)
=
\frac{\partial\mathcal L}{\partial\phi}.
\label{eq:field-el}
\end{equation}
Для \eqref{eq:scalar-lagrangian},
\begin{equation}
\frac{\partial\mathcal L_\phi}{\partial\phi_{,\mu}}
=
-\partial^\mu\phi,
\qquad
\frac{\partial\mathcal L_\phi}{\partial\phi}
=
-U'(\phi),
\end{equation}
и, следовательно,
\begin{equation}
\Box\phi=U'(\phi),
\qquad
\Box\equiv\partial_\mu\partial^\mu.
\end{equation}
Эквивалентно,
\begin{equation}
\ddot\phi-\nabla^2\phi=-U'(\phi).
\label{eq:scalar-wave}
\end{equation}
На доске перед тем, как члены с \(y\) и \(z\) оставляются неявными, явно записана только производная по \(x\). Результат — знакомое волновое уравнение
с силовым членом от потенциала.

\begin{figure}[H]
\centering
\includegraphics[width=\linewidth]{lecture_09_figure_09.png}
\caption{Локальность, лоренц-инвариантный скалярный лагранжиан и
уравнение Эйлера--Лагранжа в теории поля.}
\par\smallskip
\textit{Кадр лекции: 00:41:00.}
\end{figure}

Электромагнитное вычисление — тот же вариационный процесс, но с
четырьмя полями и большим числом индексов.

\section{Калибровочная инвариантность выбирает максвелловский скаляр}

Добавим третий принцип: калибровочную инвариантность. Пока положим
\(J^\mu=0\) и рассмотрим свободное поле. Калибровочное преобразование имеет вид
\begin{equation}
A_\mu\longrightarrow A_\mu+\partial_\mu\chi.
\label{eq:gauge-transform}
\end{equation}
Тензор поля не меняется, потому что смешанные вторые производные
коммутируют:
\begin{align}
F_{\mu\nu}
&\longrightarrow
F_{\mu\nu}
+\partial_\mu\partial_\nu\chi
-\partial_\nu\partial_\mu\chi
\nonumber\\
&=F_{\mu\nu}.
\end{align}

Одна лишь лоренц-инвариантность допускала бы скаляр
\begin{equation}
A_\mu A^\mu.
\end{equation}
Тем не менее в теории Максвелла он запрещён, потому что меняется при
\eqref{eq:gauge-transform}. Индексы в этом выражении обозначают компоненты
поля; это не производные. Производная обозначается запятой, как в
\(A_{\nu,\mu}=\partial_\mu A_\nu\).

\begin{figure}[H]
\centering
\includegraphics[width=\linewidth]{lecture_09_figure_10.png}
\caption{Калибровочное преобразование и лоренц-скаляр
\(A_\mu A^\mu\), который отвергается, потому что не является калибровочно-инвариантным.}
\par\smallskip
\textit{Кадр лекции: 00:43:40.}
\end{figure}

Следовательно, свободный лагранжиан следует строить из \(F_{\mu\nu}\).
Линейный след не проходит:
\begin{equation}
F^\mu{}_\mu=0.
\end{equation}
Антисимметрия заставляет каждую диагональную компоненту обращаться в нуль:
\begin{equation}
F_{00}=F_{11}=F_{22}=F_{33}=0.
\end{equation}
Поднятие индекса может изменить знак, но не может превратить нуль во что-то другое.

\begin{figure}[H]
\centering
\includegraphics[width=\linewidth]{lecture_09_figure_11.png}
\caption{Попытка взять линейный след и обращающиеся в нуль диагональные элементы
антисимметричного тензора поля.}
\par\smallskip
\textit{Кадр лекции: 00:47:00.}
\end{figure}

Первый ненулевой лоренц-скаляр квадратичен:
\begin{equation}
F_{\mu\nu}F^{\mu\nu}.
\label{eq:f2-scalar}
\end{equation}
Давайте распишем свёртку, а не просто приведём готовый результат. Смешанные компоненты
содержат электрическое поле. Поднятие их временного индекса даёт
знак минус, и встречаются оба порядка \(0i\) и \(i0\):
\begin{equation}
\sum_i\left(F_{0i}F^{0i}+F_{i0}F^{i0}\right)
=-2E^2.
\end{equation}
Чисто пространственные компоненты содержат магнитное поле. Поднятие
пространственных индексов знака не меняет, и оба порядка снова встречаются:
\begin{equation}
\sum_{i<j}\left(F_{ij}F^{ij}+F_{ji}F^{ji}\right)
=2B^2.
\end{equation}
Следовательно,
\begin{equation}
F_{\mu\nu}F^{\mu\nu}=2(B^2-E^2),
\qquad
-F_{\mu\nu}F^{\mu\nu}=2(E^2-B^2).
\label{eq:f2-eb}
\end{equation}

\begin{figure}[H]
\centering
\includegraphics[width=\linewidth]{lecture_09_figure_12.png}
\caption{Полностью вычисленная свёртка, выраженная через модули электрического и магнитного
полей.}
\par\smallskip
\textit{Кадр лекции: 00:51:00.}
\end{figure}

\begin{classroomqa}
\audiencequestion{Равен ли \(F_{00}\) тоже нулю?}
\lecturerresponse{Да. Каждая диагональная компонента антисимметричного
тензора равна нулю, включая \(F_{00}\). Вклад дают только недиагональные смешанные
компоненты, которые кодируют \(\mathbf E\), и недиагональные пространственные
компоненты, которые кодируют \(\mathbf B\).
\lecturetimestamp{00:51:36}}
\end{classroomqa}

Стандартная нормировка имеет вид
\begin{equation}
\boxed{
\mathcal L_{\mathrm{EM}}
=
-\frac14F_{\mu\nu}F^{\mu\nu}
}
\label{eq:maxwell-lagrangian}
\end{equation}
или
\begin{equation}
\mathcal L_{\mathrm{EM}}
=
\frac12(E^2-B^2).
\label{eq:maxwell-lagrangian-eb}
\end{equation}
Знак минус даёт электрическому члену привычный положительный знак. Множитель
\(1/4\) компенсирует двукратное появление каждой антисимметричной
пары и оставляет стандартный множитель \(1/2\). Умножение всего
свободного действия на ненулевую константу не изменило бы его уравнений
без источников, но именно эта нормировка согласуется со стандартной
связью с источником.

\section{Варьирование свободного электромагнитного поля}

Независимыми полями являются
\begin{equation}
A_0,\qquad A_x,\qquad A_y,\qquad A_z.
\end{equation}
Для каждой компоненты \(A_\nu\) уравнение Эйлера--Лагранжа имеет вид
\begin{equation}
\partial_\mu\left(
\frac{\partial\mathcal L}{\partial A_{\nu,\mu}}
\right)
=
\frac{\partial\mathcal L}{\partial A_\nu},
\qquad
A_{\nu,\mu}\equiv\partial_\mu A_\nu.
\label{eq:vector-el-comma}
\end{equation}
Расположение индексов в обозначении с запятой стоит читать медленно: первый
индекс указывает, какая компонента \(A\) является полем; второй указывает
координату дифференцирования.

В этом обозначении,
\begin{equation}
F_{\mu\nu}
=
A_{\nu,\mu}-A_{\mu,\nu}.
\label{eq:f-comma-consistent}
\end{equation}
Это в точности \eqref{eq:course-f-convention}. Если словесно читать индексы при запятой
в противоположном порядке, меняется промежуточный знак минус, поэтому
покомпонентное вычисление полезно.

\begin{figure}[H]
\centering
\includegraphics[width=\linewidth]{lecture_09_figure_13.png}
\caption{Тензор поля, записанный в обозначении с запятой после того, как четыре
компоненты потенциала были выделены как независимые поля.}
\par\smallskip
\textit{Кадр лекции: 00:55:30.}
\end{figure}

\begin{figure}[H]
\centering
\includegraphics[width=\linewidth]{lecture_09_figure_14.png}
\caption{Лагранжиан Максвелла, явно записанный через производные
векторного потенциала.}
\par\smallskip
\textit{Кадр лекции: 00:56:45.}
\end{figure}

\subsection{Одна показательная производная}

Выберем производную \(A_{x,y}=\partial_yA_x\). Члены \(xy\) и \(yx\)
в полной сумме дают вклад
\begin{align}
\mathcal L_{\mathrm{EM}}
&\supset
-\frac14\left(F_{xy}F^{xy}+F_{yx}F^{yx}\right)
\nonumber\\
&=
-\frac12F_{xy}^2
\nonumber\\
&=
-\frac12\left(A_{x,y}-A_{y,x}\right)^2.
\label{eq:isolated-xy-piece}
\end{align}
Последнее равенство корректно, потому что при возведении в квадрат знак исчезает:
\(A_{x,y}-A_{y,x}=F_{yx}=-F_{xy}\).

\begin{figure}[H]
\centering
\includegraphics[width=\linewidth]{lecture_09_figure_15.png}
\caption{Показательная производная \(A_{x,y}\) и выделенный
квадрат антисимметричной комбинации.}
\par\smallskip
\textit{Кадр лекции: 01:00:45.}
\end{figure}

Продифференцируем:
\begin{align}
\frac{\partial\mathcal L_{\mathrm{EM}}}
{\partial A_{x,y}}
&=
-\left(A_{x,y}-A_{y,x}\right)
\nonumber\\
&=
F_{xy}.
\label{eq:component-variation}
\end{align}
На доске это суммировано так, что знак минус приписан тензору, чьи
индексы идут в порядке записи производной на доске. При используемом во всём курсе
определении \eqref{eq:course-f-convention} инвариантное утверждение имеет вид
\begin{equation}
\frac{\partial\mathcal L_{\mathrm{EM}}}
{\partial(\partial_\mu A_\nu)}
=
-F^{\mu\nu}.
\label{eq:maxwell-canonical-derivative}
\end{equation}
Для примера \(\partial_yA_x\) это даёт
\(-F^{yx}=F^{xy}\), в точности
\eqref{eq:component-variation}. Таким образом, кажущееся различие знаков
сводится лишь к вопросу учёта порядка индексов.

\begin{figure}[H]
\centering
\includegraphics[width=\linewidth]{lecture_09_figure_16.png}
\caption{Раскрытие выбранного квадрата перед дифференцированием по
одной производной потенциала.}
\par\smallskip
\textit{Кадр лекции: 01:03:00.}
\end{figure}

\begin{classroomqa}
\audiencequestion{Что произошло с множителем \(1/2\) в покомпонентной
производной?}
\lecturerresponse{При дифференцировании квадрата появляется множитель два,
который сокращает \(1/2\):
\[
\frac{d}{dx}\left[-\frac12x^2\right]=-x.
\]
Прежний переход от \(1/4\) к \(1/2\) возник из-за двух упорядоченных
тензорных компонент \(xy\) и \(yx\). \lecturetimestamp{01:03:29}}
\end{classroomqa}

В вакуумном лагранжиане потенциал без производных отсутствует:
\begin{equation}
\frac{\partial\mathcal L_{\mathrm{EM}}}{\partial A_\nu}=0.
\end{equation}
Подстановка в \eqref{eq:vector-el-comma} даёт
\begin{equation}
-\partial_\mu F^{\mu\nu}=0,
\end{equation}
или, эквивалентно,
\begin{equation}
\boxed{\partial_\mu F^{\mu\nu}=0.}
\label{eq:free-maxwell-from-action}
\end{equation}

\begin{figure}[H]
\centering
\includegraphics[width=\linewidth]{lecture_09_figure_17.png}
\caption{Покомпонентный результат в обобщённом виде, подставленный в
уравнение Эйлера--Лагранжа.}
\par\smallskip
\textit{Кадр лекции: 01:09:00.}
\end{figure}

Тот же вывод получается чище, если варьировать всё действие:
\begin{align}
\delta S_{\mathrm{EM}}
&=
-\frac12\int d^4x\,
F^{\mu\nu}
\left(\partial_\mu\delta A_\nu
-\partial_\nu\delta A_\mu\right)
\nonumber\\
&=
-\int d^4x\,F^{\mu\nu}\partial_\mu\delta A_\nu
\nonumber\\
&=
\int d^4x\,
\left(\partial_\mu F^{\mu\nu}\right)\delta A_\nu,
\label{eq:global-maxwell-variation}
\end{align}
где граничный член отбрасывается. Антисимметрия объединяет два
члена в первой строке. Поскольку \(\delta A_\nu\) произвольно,
следует \eqref{eq:free-maxwell-from-action}.

\begin{classroomqa}
\audiencequestion{Не предполагал ли выделенный квадрат, что магнитное поле
направлено только вдоль \(z\)?}
\lecturerresponse{Нет. Это был один член, вынесенный из полной суммы для
рассмотрения; после него должно стоять многоточие. Оставшиеся пространственные пары
dают другие компоненты магнитного поля, а пары время--пространство дают
электрические компоненты. \lecturetimestamp{01:08:05}}
\end{classroomqa}

\begin{classroomqa}
\audiencequestion{Изменится ли результат, если выбрать один временной индекс вместо двух пространственных индексов?}
\lecturerresponse{Получается та же общая производная, а знак метрики учитывается при поднятии индекса. Два пространственных индекса выделяют магнитную компоненту; один временной и один пространственный индекс выделяют электрическую компоненту. Антисимметричная формула \eqref{eq:maxwell-canonical-derivative} охватывает оба случая.
\lecturetimestamp{01:08:46}}
\end{classroomqa}

\begin{classroomqa}
\audiencequestion{Мы всё ещё предполагаем \(J^\mu=0\)?}
\lecturerresponse{Только для этого вывода для свободного поля. Ток был временно положен равным нулю, чтобы можно было выделить калибровочно-инвариантный кинетический член. Теперь мы восстанавливаем его и спрашиваем, какой дополнительный член даёт правую часть уравнения Максвелла. \lecturetimestamp{01:09:41}}
\end{classroomqa}

Алгебра не трудна, но её нужно проделать, а не просто наблюдать. Однократное прослеживание цепочки индексов не закрепляет её надолго; а воспроизведение этой цепочки от руки закрепляет.

\section{Сохраняющийся четырёхток}

Источником является четырёхток
\begin{equation}
J^\mu=(\rho,\mathbf J).
\end{equation}
Его временная компонента — это плотность заряда, а пространственные компоненты — потоки заряда через три взаимно перпендикулярных окна. Эквивалентно, для небольшого неподвижного ящика,
\begin{equation}
\frac{\partial\rho}{\partial t}
+\nabla\cdot\mathbf J=0.
\label{eq:continuity-vector-l9}
\end{equation}
Заряд внутри ящика может уменьшаться только за счёт протекания через его границу. В ковариантной записи,
\begin{equation}
\boxed{\partial_\mu J^\mu=0.}
\label{eq:continuity-covariant-l9}
\end{equation}

\begin{figure}[H]
\centering
\includegraphics[width=\linewidth]{lecture_09_figure_18.png}
\caption{Напоминание о плотности заряда, пространственном токе и интерпретации через малый ящик перед введением связи с источником.}
\par\smallskip
\textit{Кадр лекции: 01:14:45.}
\end{figure}

\section{Связь с источником и калибровочная инвариантность}

Простейший локальный лоренц-инвариантный скаляр, линейный как по току, так и по потенциалу, имеет вид
\begin{equation}
S_{\mathrm{int}}
=
-\int d^4x\,J^\mu A_\mu.
\label{eq:interaction-action}
\end{equation}
Общий знак задаётся соглашением. В лекции в шутку говорится, что исторический выбор знака электрического заряда в конечном счёте можно свалить на Benjamin Franklin; реальная суть в том, что всюду нужно использовать одну и ту же согласованную знаковую конвенцию.

На первый взгляд \eqref{eq:interaction-action} не выглядит калибровочно-инвариантным.
При преобразовании \eqref{eq:gauge-transform},
\begin{equation}
\delta_\chi S_{\mathrm{int}}
=
-\int d^4x\,J^\mu\partial_\mu\chi.
\label{eq:gauge-variation-interaction}
\end{equation}

\begin{figure}[H]
\centering
\includegraphics[width=\linewidth]{lecture_09_figure_19.png}
\caption{Действие взаимодействия и его вариация при \(A_\mu\mapsto A_\mu+\partial_\mu\chi\).}
\par\smallskip
\textit{Кадр лекции: 01:16:45.}
\end{figure}

Чтобы увидеть механизм в элементарных терминах, разложим свёртку:
\begin{equation}
J^\mu\partial_\mu\chi
=
J^0\partial_0\chi
+J^1\partial_1\chi
+J^2\partial_2\chi
+J^3\partial_3\chi.
\end{equation}
Предположим, что ток ограничен большой лабораторией или достаточно быстро спадает, так что на границе он обращается в нуль. Проинтегрируем по частям каждый член. Например, для первой пространственной координаты,
\begin{align}
-\int d^4x\,J^1\partial_1\chi
&=
-\int dt\,dx^2\,dx^3
\left[J^1\chi\right]_{x^1=-\infty}^{x^1=+\infty}
\nonumber\\
&\quad
+\int d^4x\,\chi\,\partial_1J^1
\nonumber\\
&=
\int d^4x\,\chi\,\partial_1J^1.
\end{align}
Та же операция применяется к временному члену и двум другим пространственным членам.

\begin{figure}[H]
\centering
\includegraphics[width=\linewidth]{lecture_09_figure_20.png}
\caption{Четыре члена калибровочной вариации, разложенные до интегрирования по частям.}
\par\smallskip
\textit{Кадр лекции: 01:19:15.}
\end{figure}

\begin{figure}[H]
\centering
\includegraphics[width=\linewidth]{lecture_09_figure_21.png}
\caption{Показательное интегрирование по частям переносит производную с калибровочной функции на ток.}
\par\smallskip
\textit{Кадр лекции: 01:20:45.}
\end{figure}

Собирая все четыре компоненты,
\begin{equation}
\delta_\chi S_{\mathrm{int}}
=
\int d^4x\,\chi\,\partial_\mu J^\mu.
\label{eq:gauge-variation-continuity}
\end{equation}
Следовательно,
\begin{equation}
\partial_\mu J^\mu=0
\quad\Longrightarrow\quad
\delta_\chi S_{\mathrm{int}}=0.
\end{equation}
При принятых граничных условиях взаимодействие с источником является калибровочно-инвариантным в точности тогда, когда ток локально сохраняется.

\begin{classroomqa}
\audiencequestion{Достаточно ли потребовать только \(\nabla\cdot\mathbf J=0\)?}
\lecturerresponse{Нет. Нужны все четыре члена:
\[
\partial_\mu J^\mu
=
\dot\rho+\nabla\cdot\mathbf J=0.
\]
Первый член — это изменение плотности заряда в неподвижном ящике. Ненулевая пространственная дивергенция допустима, если она уравновешивается этим изменением. \lecturetimestamp{01:22:00}}
\end{classroomqa}

\section{Уравнение Максвелла с источником}

Полная электромагнитная плотность лагранжиана имеет вид
\begin{equation}
\boxed{
\mathcal L
=
-\frac14F_{\mu\nu}F^{\mu\nu}
-J^\mu A_\mu.
}
\label{eq:full-maxwell-lagrangian}
\end{equation}
Член источника не содержит производной от \(A_\nu\), поэтому он оставляет \eqref{eq:maxwell-canonical-derivative} без изменения и даёт вклад
\begin{equation}
\frac{\partial\mathcal L}{\partial A_\nu}=-J^\nu.
\end{equation}
Поэтому уравнение Эйлера--Лагранжа даёт
\begin{equation}
-\partial_\mu F^{\mu\nu}=-J^\nu,
\end{equation}
или
\begin{equation}
\boxed{\partial_\mu F^{\mu\nu}=J^\nu.}
\label{eq:sourced-maxwell-from-action}
\end{equation}

\begin{figure}[H]
\centering
\includegraphics[width=\linewidth]{lecture_09_figure_22.png}
\caption{Свободный член и член источника, показанные вместе с характерными производными Эйлера--Лагранжа.}
\par\smallskip
\textit{Кадр лекции: 01:25:00.}
\end{figure}

Несохраняющийся заданный ток привёл бы сразу к двум сбоям. Действие взаимодействия не было бы калибровочно-инвариантным, а уравнения Максвелла с источником были бы несовместны. Заряд не может просто исчезнуть внутри области без компенсирующего потока через её границу.

\section{Уравнения Максвелла снова дают закон сохранения}

Совместимость можно прочитать прямо из уравнения поля. Применяя \(\partial_\nu\) к \eqref{eq:sourced-maxwell-from-action}, получаем
\begin{equation}
\partial_\nu J^\nu
=
\partial_\nu\partial_\mu F^{\mu\nu}
=0,
\label{eq:antisymmetry-implies-continuity}
\end{equation}
поскольку производные симметричны относительно \(\mu\leftrightarrow\nu\), тогда как \(F^{\mu\nu}\) антисимметричен.

В лекции это также проверяется в трёхвекторной форме, и эта проверка выявляет сделанную по ходу дела ошибку в знаке. Начнём с закона Гаусса:
\begin{equation}
\rho=\nabla\cdot\mathbf E.
\end{equation}
Продифференцируем по времени:
\begin{equation}
\dot\rho=\nabla\cdot\dot{\mathbf E}.
\label{eq:rho-dot-div-edot}
\end{equation}
На доске предварительно было записано уравнение
\begin{equation}
\dot{\mathbf E}
=
-\nabla\times\mathbf B+\mathbf J.
\label{eq:provisional-wrong-source-sign}
\end{equation}
Его дивергенция дала бы
\begin{equation}
\dot\rho=\nabla\cdot\mathbf J,
\end{equation}
поскольку дивергенция ротора равна нулю. Это противоречит
\eqref{eq:continuity-vector-l9}.

\begin{figure}[H]
\centering
\includegraphics[width=\linewidth]{lecture_09_figure_23.png}
\caption{Предварительное уравнение на доске с \(\mathbf J\) со знаком плюс, дающее неверный знак в уравнении непрерывности.}
\par\smallskip
\textit{Кадр лекции: 01:27:00.}
\end{figure}

Несоответствие не скрывается. Вычисление заново прослеживается в реальном времени; формулируется требуемое уравнение непрерывности; и один из слушателей проверяет уравнение и указывает знак при источнике. Согласованное с обычной конвенцией компонентное уравнение имеет вид
\begin{equation}
\boxed{
\dot{\mathbf E}+\nabla\times\mathbf B=-\,\mathbf J,
}
\label{eq:corrected-source-sign}
\end{equation}
или
\begin{equation}
\dot{\mathbf E}
=
-\nabla\times\mathbf B-\mathbf J.
\end{equation}
Теперь уравнения \eqref{eq:rho-dot-div-edot} и
\eqref{eq:corrected-source-sign} дают
\begin{align}
\dot\rho
&=
\nabla\cdot
\left(-\nabla\times\mathbf B-\mathbf J\right)
\nonumber\\
&=
-\nabla\cdot\mathbf J,
\end{align}
что в точности и есть локальное сохранение заряда.

\begin{figure}[H]
\centering
\includegraphics[width=\linewidth]{lecture_09_figure_24.png}
\caption{Исправленный знак при источнике и восстановленное уравнение непрерывности.}
\par\smallskip
\textit{Кадр лекции: 01:31:00.}
\end{figure}

\begin{classroomqa}
\audiencequestion{Должен ли член с током справа быть \(-\mathbf J\)?}
\lecturerresponse{Да. При соглашении для тензора поля, используемом на протяжении всего курса, пространственная часть \(\partial_\mu F^{\mu\nu}=J^\nu\) имеет вид \(\dot{\mathbf E}+\nabla\times\mathbf B=-\mathbf J\).
Это исправление приводит к тому, что вычисление дивергенции воспроизводит \(\dot\rho+\nabla\cdot\mathbf J=0\), и не изменяет предыдущего ковариантного вывода из действия. \lecturetimestamp{01:30:17}}
\end{classroomqa}

\section{Какой принцип является фундаментальным?}

\begin{classroomqa}
\audiencequestion{Является ли это рассуждение круговым? Мы использовали сохранение заряда, чтобы сделать действие калибровочно-инвариантным, затем использовали действие для вывода уравнений Максвелла, а уравнения Максвелла влекут сохранение заряда.}
\lecturerresponse{Эти утверждения взаимно совместимы, а не образуют круговое доказательство, предполагающее собственный вывод. Можно начать с наблюдаемых уравнений Максвелла и вывести уравнение непрерывности; тогда \(-J^\mu A_\mu\) распознаётся как калибровочно-инвариантная связь. Или можно начать с локальности, лоренц-инвариантности, калибровочной инвариантности и действия; калибровочная инвариантность требует сохраняющегося тока, а действие даёт уравнения Максвелла. Или можно считать первичным физический факт, что заряд не исчезает локально. Каждый путь приводит к одной и той же структуре. \lecturetimestamp{01:31:11}}
\end{classroomqa}

Если спрашивать о наиболее экономном конструктивном пути, то сочетание локальности, лоренц-инвариантности, калибровочной инвариантности и лагранжева принципа выделяет
\begin{equation}
-\frac14F_{\mu\nu}F^{\mu\nu}-J^\mu A_\mu
\end{equation}
и тем самым динамическое уравнение Максвелла. Называть один элемент этого треугольника более фундаментальным — отчасти вопрос вкуса. Важный факт состоит в том, что эти элементы согласуются.

\section{Почему слово \emph{калибровка}?}

\begin{classroomqa}
\audiencequestion{Математика ясна, но что здесь означает слово \emph{калибровка}?}
\lecturerresponse{Это название — исторический анахронизм. Калибровка может относиться к выбранному нулю, как, например, давление, измеряемое относительно некоторого эталона. Hermann Weyl первоначально связал эту идею с локальным эталоном длины; такая интерпретация не была правильной физикой. Термин сохранился, когда эту свободу правильно связали с фазой квантовой волновой функции. В данном классическом вычислении замена \(A_\mu\) на \(\partial_\mu\chi\) меняет условный нуль описания, не изменяя \(F_{\mu\nu}\). Более глубокий квантовый смысл пока откладывается.
\lecturetimestamp{01:34:49}}
\end{classroomqa}

Исторический ярлык не должен заслонять рабочее утверждение: потенциалы, связанные преобразованием \eqref{eq:gauge-transform}, описывают одно и то же электромагнитное поле. Поэтому калибровочная свобода является избыточностью описания через потенциал, тогда как тензор поля калибровочно-инвариантен.

\section{Почему два семейства уравнений Максвелла выглядят несимметрично}

\begin{classroomqa}
\audiencequestion{Если электрическое и магнитное поля столь симметричны, почему одно из семейств уравнений Максвелла выступает как тождество Бианки, тогда как другое получается из уравнений движения?}
\lecturerresponse{В электромагнетизме без источников можно ввести дуальный потенциал так, чтобы роли поменялись местами: то, что было уравнением движения, становится тождеством Бианки, и наоборот. Физическая асимметрия вносится источниками. Электрические заряды наблюдаются, тогда как элементарные магнитные монополи не были обнаружены. Если бы существовала точная обычная симметрия, переставляющая электрические и магнитные заряды, следовало бы ожидать магнитных аналогов с соответствующими спектрами, но их не наблюдают. Более сложные теории всё же могут содержать магнитные заряды или электрические--магнитные дуальности, но это выходит за рамки данного вывода. \lecturetimestamp{01:37:11}}
\end{classroomqa}

Таким образом, уравнения без источников обладают структурой, допускающей электрическо--магнитную перестановку, тогда как наблюдаемый сектор источников выбирает электрическое описание. Это больше, чем типографская условность, но это не устраняет симметрию, видимую в свободной плоской волне.

\section{Что было собрано воедино}

Теперь путь по лекции завершён:
\begin{enumerate}
\item Уравнения Максвелла заставляют плоскую волну быть поперечной, с ортогональными электрическим и магнитным полями и дисперсионным соотношением \(\omega=ck\).
\item Уравнения Бианки — это тождества, следующие из потенциала; свёрнутые уравнения являются динамическими.
\item Локальность и инвариантность Лоренца приводят нас к плотности действия, тогда как калибровочная инвариантность исключает \(A_\mu A^\mu\) и выделяет скаляр, построенный из \(F_{\mu\nu}\).
\item Свободное действие \(-\tfrac14F_{\mu\nu}F^{\mu\nu}\) даёт вакуумное уравнение Максвелла.
\item Взаимодействие \(-J^\mu A_\mu\) является калибровочно-инвариантным в точности тогда, когда ток удовлетворяет уравнению непрерывности.
\item Действие с источником даёт \(\partial_\mu F^{\mu\nu}=J^\nu\), а из антисимметрии обратно получается \(\partial_\mu J^\mu=0\).
\end{enumerate}

Это очень богатая структура: лагранжева механика, релятивистская теория поля, электродинамика, инвариантность Лоренца, калибровочная инвариантность и локальное сохранение заряда стали частями единого рассуждения. Единственный надёжный способ это усвоить — снова начать с уравнений и вручную проследить каждую производную и каждый знак.
